\section{Introducción: de Flow Matching 1 a Flow Matching 2}

Esta clase es la número 10 del curso KAIST CS492D (Modelos de difusión y modelos de flujo) y constituye la segunda parte del módulo sobre Flow Matching. El ponente Minhyuk Sung primero repasa el marco básico de Flow Matching introducido en la clase anterior —los tres conceptos centrales de mapa de flujo (Flow Map), campo vectorial (Vector Field) y trayectoria de probabilidad (Probability Path)—, sus relaciones mutuas y, a partir de esta base, profundiza en los siguientes temas clave:

\begin{enumerate}
    \item La \textbf{equivalencia matemática} entre modelos de difusión y modelos de flujo: bajo qué condiciones son equivalentes y cómo se transforman uno al otro?
    \item La \textbf{correspondencia entre campo de velocidad y función score}: ¿cuál es la expresión matemática exacta entre velocity y score function?
    \item \textbf{Rectified Flow}: ¿cómo simplifica el proceso de generación las trayectorias rectas y cómo permite el ajuste fino con Reflow hacerlas aún más lineales?
    \item Perspectiva de \textbf{transporte óptimo}: ¿cómo afectan los emparejamientos aleatorios frente a los óptimos al grado de curvatura de las trayectorias?
    \item \textbf{Aplicaciones prácticas}: ¿cómo aplican las tecnologías de Flow Matching los modelos Stable Diffusion 3 (SD3) y Flux?
\end{enumerate}

\begin{knowledgebox}{Repaso de conocimientos previos}
Para comprender esta clase se requieren los siguientes conceptos previos introducidos en la clase anterior (Lecture 9: Flow Matching 1):
\begin{itemize}
    \item \textbf{Mapa de flujo (Flow Map)} $\psi_t$: mapeo determinista desde $t=0$ hasta $t$, con $\psi_0(\mathbf{x}) = \mathbf{x}$
    \item \textbf{Campo vectorial (Vector Field)} $u_t(\mathbf{x})$: derivada temporal del mapa de flujo, $\frac{d\psi_t}{dt} = u_t(\psi_t(\mathbf{x}))$
    \item \textbf{Trayectoria de probabilidad (Probability Path)} $p_t$: distribución de probabilidad que varía con el tiempo $t$, satisfaciendo $p_0 = p_{\text{ref}}$ y $p_1 = p_{\text{data}}$
    \item \textbf{Ecuación de Fokker--Planck}: ecuación que vincula el campo vectorial con la trayectoria de probabilidad
    \item \textbf{Flow Matching condicional (Conditional Flow Matching)}: descomposición del campo vectorial marginal en una expectativa sobre campos vectoriales condicionales
\end{itemize}
\end{knowledgebox}

\subsection{Convenciones de notación y cambio del dominio temporal}

En el contexto de Flow Matching, el dominio temporal difiere del de los modelos de difusión:

\begin{itemize}
    \item \textbf{Modelos de difusión}: el tiempo va desde $T$ (extremo ruidoso) hasta $0$ (extremo de datos); la generación es un ``proceso inverso''
    \item \textbf{Modelos de flujo}: el tiempo va desde $0$ (distribución de referencia) hasta $1$ (distribución de datos); la generación es un ``proceso directo''
\end{itemize}

Específicamente:
\begin{itemize}
    \item $\mathbf{x}_0 \sim p_0 = \mathcal{N}(\mathbf{0}, \mathbf{I})$: muestra proveniente de la distribución de referencia (gausiana estándar)
    \item $\mathbf{x}_1 \sim p_1 = p_{\text{data}}$: muestra proveniente de la distribución de datos
\end{itemize}

\begin{warningbox}{Símbolos propensos a confusión}
Observe que en Flow Matching, $\mathbf{x}_0$ es el \textbf{ruido} y $\mathbf{x}_1$ son los \textbf{datos}, lo cual es exactamente opuesto a la convención habitual en DDPM donde $\mathbf{x}_0$ son datos y $\mathbf{x}_T$ es ruido. Además, las funciones $\mu_t, \sigma_t$ de la Lecture 9 son esencialmente idénticas a los $\alpha_t, \sigma_t$ de los modelos de difusión variacional; la única diferencia radica en el cambio del dominio temporal, que invierte sus posiciones relativas.
\end{warningbox}

\subsection{Forma general del mapa de flujo condicional}

En la clase anterior se introdujo la forma lineal gaussiana del mapa de flujo condicional (Conditional Flow Map):
$$\psi_t(\mathbf{x}_0 \mid \mathbf{x}_1) = \sigma_t \cdot \mathbf{x}_0 + \alpha_t \cdot \mathbf{x}_1$$

Donde:
\begin{itemize}
    \item $\sigma_t, \alpha_t$ son funciones escalares del tiempo $t$
    \item Condiciones de frontera: $\sigma_0 = 1, \alpha_0 = 0$ (en $t=0$ se obtiene puro ruido $\mathbf{x}_0$); $\sigma_1 = 0, \alpha_1 = 1$ (en $t=1$ se obtienen datos $\mathbf{x}_1$)
\end{itemize}

La trayectoria de probabilidad condicional correspondiente es una distribución gaussiana:
$$p_t(\mathbf{x} \mid \mathbf{x}_1) = \mathcal{N}(\mathbf{x}; \alpha_t \mathbf{x}_1, \sigma_t^2 \mathbf{I})$$

Tomando la derivada temporal del mapa de flujo condicional se obtiene el campo vectorial condicional:
$$u_t(\mathbf{x} \mid \mathbf{x}_1) = \dot{\sigma}_t \cdot \mathbf{x}_0 + \dot{\alpha}_t \cdot \mathbf{x}_1$$

\subsection{Resumen de la sección}

En esta sección se repasan los supuestos básicos de Flow Matching, incluyendo las convenciones del dominio temporal, la definición de mapa de flujo condicional y la derivación del campo vectorial condicional. El punto clave es que en el modelo de flujo el tiempo va desde $0$ (ruido) hasta $1$ (datos), y el objetivo del entrenamiento es aprender una red neuronal para predecir el campo vectorial (velocidad); esto se asemeja conceptualmente a la predicción de ruido en los modelos de difusión.

